% year: תשפ"ג
% type: מבחן
% semester: א
% moed: ב
% number: 6א
% points: 10
% categories: integral-applications
% source: exams/מבחנים/תשפג/מועד ב תשפג.docx

\begin{question}
(10 נק') חשבו את השטח הסופי החסום בין גרף הפונקציה $f(x)=x^3+x^2-2x$ וציר ה-$x$.
\end{question}

\begin{hint}
פרקו את הפולינום לגורמים כדי למצוא את נקודות החיתוך עם ציר $x$ ואת הסימן בכל קטע.
\end{hint}

\begin{solution}
פירוק: $f(x)=x(x^2+x-2)=x(x+2)(x-1)$. נקודות החיתוך עם ציר $x$: $x=-2,0,1$. סימן: ב-$(-2,0)$ $f>0$ (למשל $f(-1)=2$), וב-$(0,1)$ $f<0$ (למשל $f(\frac12)=-\frac58$). התחום הסופי החסום בין הגרף לציר הוא מעל $[-2,1]$.

פונקציה קדומה: $F(x)=\frac{x^4}{4}+\frac{x^3}{3}-x^2$. אז $F(-2)=4-\frac83-4=-\frac83$, $F(0)=0$, $F(1)=\frac14+\frac13-1=-\frac{5}{12}$.
\[ S=\int_{-2}^0 f\,dx-\int_0^1 f\,dx=\big(F(0)-F(-2)\big)-\big(F(1)-F(0)\big)=\frac83+\frac{5}{12}=\frac{37}{12}. \]
\textbf{תשובה:} $S=\frac{37}{12}$.
\end{solution}
